\section{Conventions and companion theorems}\label{sec:preliminaries}

We assume familiarity with discrepancies, singularities of pairs, and
the minimal model program, as in \cite{KM,KollarSing}. We recall the
conventions and precise theorem inputs needed for the argument.
The argument is carried out over $\C$ until the final descent to an
arbitrary algebraically closed field of characteristic zero. Varieties
are integral and projective unless a different setting is specified.
A \emph{contraction} is a projective surjective morphism
$f:X\to Z$ of normal varieties satisfying $f_*\OO_X=\OO_Z$.
In characteristic zero its generic fibre is geometrically integral.
All boundaries in our argument are rational.

For a normal variety, a rational section of a divisorial sheaf is
identified with a rational function. Thus, for a rational Weil divisor $D$,
\begin{equation}\label{eq:section-inequality}
 H^0(X,\OO_X(\floor{D}))
 =\{a\in k(X):\Div(a)+D\geq0\}\cup\{0\}.
\end{equation}
The equality uses the integrality of the orders of $a$. It does not
require $\floor{D}$ to be Cartier. We use compatible canonical divisors
on birational models, obtained from one rational top differential.
A rational $m$-canonical form means an element of the $m$th tensor power
of the top differential line of a function field; its divisor is computed
on normal models at codimension-one regular points.

We use log discrepancies. If $\pi:Y\to X$ is a birational model and
$K_Y+B_Y=\pi^*(K_X+B)$, then the log discrepancy of a prime divisor $E$
on $Y$ is $1-\operatorname{coeff}_E B_Y$. In particular, log canonicity
means that all these numbers are nonnegative. Crepant boundaries on
higher models may have negative coefficients. Effectiveness will always
be specified when it is needed. We use dlt adjunction in its usual form:
a component $S$ of the coefficient-one part of a dlt boundary is normal,
and
\[
 (K_Y+B_Y)|_S=K_S+\operatorname{Diff}_S(B_Y-S)
\]
with an effective lc different when $B_Y$ is an effective boundary
\cite{KM,KollarSing}.

\subsection{The three companion inputs}

The following statements are the substantive inputs from companion
manuscripts. We recall their full scope because the coefficient and
ground-field conditions matter at different places in the proof.

\begin{theorem}[Normal log canonical indices {\cite[Theorem~1.2]{U}}]
\label{thm:normal-index}
For every integer $n\geq0$ and every DCC set
$\Psi\subset[0,1]\cap\Q$, there is an integer $a(n,\Psi)>0$
such that every normal integral projective lc pair $(V,\Delta)$ over
$\C$ with $\dim V=n$, $\Delta\geq0$, coefficients in $\Psi$, and
$K_V+\Delta\simq0$ satisfies
\[
                     a(n,\Psi)(K_V+\Delta)\sim0.
\]
Here the displayed multiple is an integral principal divisor.
\end{theorem}

A DCC set is a set with no infinite strictly decreasing sequence.
We will apply Theorem~\ref{thm:normal-index} both to a finite coefficient
set and to the DCC set produced by finite cyclic quotients. Its
principal-divisor conclusion supplies an actual rational pluriform
in a bounded degree.

\begin{theorem}[Good models {\cite[Theorem~11.1]{LA}}]
\label{thm:good-models}
Every projective lc pair $(V,\Delta)$ over $\C$ with effective
real boundary and pseudo-effective real Cartier adjoint
$K_V+\Delta$ has a good log minimal model.
\end{theorem}

We use only the rational-boundary case, including semiampleness of a nef
adjoint. This theorem supplies existence; the terminating programs used
below also use the MMP with scaling \cite{TX}. Section~\ref{sec:mmp}
explains precisely how to apply these results over a projective curve.

\begin{theorem}[Arithmetic Stein degrees {\cite[Main Theorem]{SD}}]
\label{thm:stein-degree}
For every integer $n\geq1$ and every real number $t>0$, there is
an integer $N(n,t)$ with the following property. Let $F$ be a field of
characteristic zero, and let $(V,\Delta)$ be a normal integral projective
lc rational pair over $F$ with
\[
 \dim V=n,\qquad H^0(V,\OO_V)=F,\qquad
 \Delta\geq0,\qquad K_V+\Delta\simq0.
\]
If $S$ is a prime component of $\Delta$ of coefficient at least $t$,
the algebraic closure of $F$ in $F(S)$ has degree at most $N(n,t)$
over $F$.
\end{theorem}

For a proper normal component this algebraic closure is its field of
global functions. In our application $F$ is a complex curve field and
$t=1$; the theorem controls the finite curve in the Stein factorization
of a horizontal coefficient-one component.

\subsection{Trivializations and constants}

Two elementary facts will be useful repeatedly. First, a rational
function with zero divisor on a normal integral proper variety is an
invertible global function. Second, the degree of a geometric
trivialization does not increase on descending the constants.

\begin{lemma}[Descent in the same degree]\label{lem:trivialization-descent}
Let $V$ be a geometrically integral normal projective variety over a
characteristic-zero field $F$. Let $D$ be an integral Weil divisor.
If $D_{\overline F}$ is principal, then $D$ is principal.
In particular, an integral principal multiple of a rational log
canonical divisor on $V_{\overline F}$ descends in the same degree to $V$.
\end{lemma}
\begin{proof}
The divisorial sheaf $\OO_V(D)$ becomes the trivial line bundle after
extension to $\overline F$. Divisorial sheaves commute with this extension,
as can be seen on a regular big open and by reflexive extension.
Faithfully flat descent makes $\OO_V(D)$ invertible.
Proper base change shows that its space of sections is one-dimensional
over $F$. The evaluation map from this space tensored with $\OO_V$
is an isomorphism after scalar extension, hence already an isomorphism.
A nonzero section therefore trivializes the line bundle and principalizes
$D$.
\end{proof}

The geometric instances of Theorem~\ref{thm:normal-index} over a complex
function field are legitimate: the algebraic closure of any finitely
generated extension of $\C$ is abstractly isomorphic to $\C$.
Transporting the variety and its divisor through such an isomorphism
preserves the stated algebraic hypotheses. The final section gives
a separate descent argument for the ground field in the main theorem.
